%% kapa: exam=1 round=26 year=2015 period=2 form=A subject=accounting subject_ko=회계학 seq=3 number=43
%% source: Q-Net 감정평가사 기출문제 (cst00309) artlSeq=5138001 file=extracted/1008/5138001/제26회_감정평가사 1차_기출/제26회_감정평가사 1차_기출/2교시_감평_A형.pdf
%% conversion: kordoc 4.24.51 -> tools/gen1.py (auto); review=proofread-sonnet; flags=none
\begin{question}{43}
(주)감평은 보유중인 유형자산을 (주)대한의 유형자산과 교환하면서 공정가치 차액에 해당하는 현금 ￦300,000을 지급하였다. 교환일 현재 보유중인 유형자산의 취득원가는 ￦2,100,000, 감가상각누계액은 ￦500,000, 공정가치는 ￦1,700,000이다. (주)감평이 교환과정에서 인식할 유형자산의 취득원가와 유형자산처분손익은? (단, 동 교환거래는 상업적 실질이 있다고 가정한다.)

\noindent\hspace*{2em}\makebox[8em][l]{\uline{취득원가}}\uline{유형자산처분손익}\par
\begin{choices}
  \item \makebox[8em][l]{￦2,000,000}￦0
  \item \makebox[8em][l]{￦2,000,000}이익 ￦100,000
  \item \makebox[8em][l]{￦1,900,000}손실 ￦100,000
  \item \makebox[8em][l]{￦1,900,000}이익 ￦100,000
  \item \makebox[8em][l]{￦1,900,000}￦0
\end{choices}
\end{question}
\answer{2}
